\section{终端怎样区分前台和后台作业}

键盘只产生扫描码，字符设备只搬运字节。把 Ctrl-Z 解释为“停止当前前台
管线”，需要 TTY 行规程、控制终端、session、process group、signal、
scheduler 和 Shell 作业表共同维护一条所有权链。

\topic{从电传打字机到控制终端}
终端这个词早于今天的图形窗口。电传打字机把按键编码为串行字符，同时在纸上
打印远端返回的字符；后来 CRT 终端仍保留“输入流加输出流”的接口。早期 Unix
让每个登录 Shell 关联一台终端，程序继承标准输入输出，于是出现一个不可回避
的问题：同一次登录中既有 Shell，又有 Shell 启动的许多程序，输入在一个时刻
究竟属于谁？

若只依赖父子关系，答案并不正确。一条
\code{producer | filter | consumer} 管线的三个进程互为兄弟，却应一起收到
Ctrl-C；Shell 是它们的父进程，却必须继续存活。Unix 因而把身份分成三层：

\begin{table}[htbp]
\centering
\small
\begin{osgridtabularx}{\textwidth}{l l X}
关系 & 公开身份 & 解决的问题\\ \hline
进程树 & PID 与 parent PID & 谁 wait、谁收养 Zombie\\ \hline
进程组 & PGID & 哪些兄弟进程组成一个作业和信号目标\\ \hline
会话 & SID & 哪些进程组共享一台控制终端\\ \hline
\end{osgridtabularx}
\caption{PID、PGID 与 SID 的职责不能互相替代}
\end{table}

session 通常由登录 Shell 的祖先建立；一个 session 包含多个 process group，
终端只记录其中一个 foreground process group。Shell 建立作业、切换前台组，
TTY 把控制字符定向给该组。设备驱动不解析命令，Shell 也不在前台程序运行时
偷偷读取同一输入。

\topic{整条控制流跨越哪些硬件和软件层}\par
\begin{pdfdiagram}[node distance=9mm and 11mm]
  \node[os hardware] (key) {PS/2 键盘\\Set 1 make/break};
  \node[os hardware,right=of key] (irq) {i8042\\IRQ1};
  \node[os state,right=of irq] (decode) {Ctrl 状态\\C0 字符};
  \node[os state,right=of decode] (tty) {Terminal\\行规程与前台 PGID};
  \node[os state,below=of tty] (signal) {组 signal\\INT/TSTP/CONT};
  \node[os state,left=of signal] (sched) {Scheduler\\Stopped/Ready};
  \node[os state,left=of sched] (wait) {ProcessTree\\wait event};
  \node[os block,left=of wait] (shell) {Ring 3 Shell\\jobs/fg/bg};
  \draw[os arrow] (key) -- (irq);
  \draw[os arrow] (irq) -- (decode);
  \draw[os arrow] (decode) -- (tty);
  \draw[os arrow] (tty) -- (signal);
  \draw[os arrow] (signal) -- (sched);
  \draw[os arrow] (sched) -- (wait);
  \draw[os arrow] (wait) -- (shell);
\end{pdfdiagram}
\diagramnote{Ctrl-Z 从硬件按键到 Shell 停止通知至少跨过七个可独立验证的边界。}

QEMU 只模拟第一行硬件。QMP \code{send-key ctrl-z} 不会调用 Kernel 函数，也
不会直接写 ASCII。自研固件、Stage 1、Kernel、i8042、IRQ、TTY、信号和用户
Shell 仍由目标代码完成。

\topic{Set 1 扫描码为何需要修饰键状态}
PS/2 键盘发送的是按键位置事件。左 Ctrl 的 make code 是 \code{0x1D}，
break code 把最高位置一，成为 \code{0x9D}；右 Ctrl 使用扩展前缀
\code{E0 1D/E0 9D}。字母 Z 的 make code 本身永远不会说明 Ctrl 是否仍按下，
解码器必须跨中断保存左右修饰状态：

\begin{lstlisting}
left Ctrl make       -> left_control_pressed = true
left Ctrl break      -> left_control_pressed = false
E0 right Ctrl make   -> right_control_pressed = true
E0 right Ctrl break  -> right_control_pressed = false
\end{lstlisting}

传统 ASCII 的 C0 区把 Ctrl 与字母的组合编码为低五位：

\[
  \text{Ctrl-C}=\texttt{'C'}\mathbin{\&}\texttt{0x1F}=\texttt{0x03}
\]
\[
  \text{Ctrl-D}=\texttt{'D'}\mathbin{\&}\texttt{0x1F}=\texttt{0x04},
  \qquad
  \text{Ctrl-Z}=\texttt{'Z'}\mathbin{\&}\texttt{0x1F}=\texttt{0x1A}
\]

这不是项目随意挑选的数字，而是电传终端兼容历史留给字符流接口的约定。
\filepath{device/device\_model.cpp} 只把扫描码变成字符；它不知道前台进程组，
也不直接发送信号。

\topic{行规程为何位于键盘和文件描述符之间}
如果每个可打印键都立即让阻塞 \code{read} 返回，退格会成为用户数据，用户
程序必须重复实现行编辑。canonical 行规程先把半行保存在 edit buffer，
换行后才把一整行转移到 input ring：

\begin{lstlisting}
edit:   h e l x
              Backspace
edit:   h e l
              l o Enter
input:  h e l l o \n
\end{lstlisting}

\code{Terminal} 固定拥有 512 字节编辑区、1024 字节可读输入环与
2048 字节输出环。固定容量使最坏内存、失败语义和测试边界均可推导；这不是
宣称现实终端永远只需这些字节。

\begin{table}[htbp]
\centering
\small
\begin{osgridtabularx}{\textwidth}{l X l}
输入 & 行规程动作 & 用户是否读取该控制字节\\ \hline
普通字符 & 追加编辑区 & 换行提交后读取\\ \hline
CR/LF & 提交编辑区并统一附加 LF & LF 可见\\ \hline
Backspace/Delete & 删除最后一个编辑字节 & 否\\ \hline
Ctrl-D & 提交当前内容，或登记一次 EOF & 否\\ \hline
Ctrl-C & 清除未提交行，产生 InterruptForeground & 否\\ \hline
Ctrl-Z & 清除未提交行，产生 StopForeground & 否\\ \hline
\end{osgridtabularx}
\caption{canonical 输入不是普通字节 FIFO}
\end{table}

\code{SubmitCharacter} 返回强类型 \code{TerminalInputAction}。IRQ 上层只按
Buffered、InputReady、EndOfFileReady、InterruptForeground 或 StopForeground
分派，不重复认识 \code{0x03}/\code{0x1A}。这把硬件编码与作业策略隔开。

\topic{输入守恒为何必须增加 consumed}
早期 ConsoleInput 可以要求 submitted 等于 read 加 buffered，因为每个接受
字符最终都交给用户。行规程会故意吞掉控制字符，还会撤销已经进入编辑区的
数据。一次有效退格消耗两个已提交事实：退格按键本身和被删除字符。

终端使用精确等式：

\[
S=R+B+E+D+C
\]

\begin{itemize}
  \item \(S\)：硬件路径接受的 submitted bytes；
  \item \(R\)：已经由用户读取的 bytes；
  \item \(B\)：已提交输入环中尚未读取的 bytes；
  \item \(E\)：尚未提交的 editing bytes；
  \item \(D\)：容量边界拒绝的 dropped bytes；
  \item \(C\)：控制语义消费的 control/erased bytes。
\end{itemize}

Ctrl-C 或 Ctrl-Z 清空半行时，\code{consumed} 同时增加控制字节和被撤销的
编辑长度；空编辑区的退格仍消费退格本身。这个账本不仅让整机终态可验证，也
能定位“控制键后丢了一个普通字符”到底是预期撤销还是环索引损坏。

输出侧满足：

\[
Q=W+P
\]

\(Q\) 是 queued，\(W\) 是 written，\(P\) 是 pending。阶段结束要求
\(P=0\)，所以所有成功接收的目标输出都已经到达串口设备。

\topic{单核为什么仍需要 irq-save 锁}
单 BSP 只排除了另一个 CPU 并行执行，没有排除同核中断重入。Thread 在系统
调用中修改输出环时，IRQ1 可以打断它并提交输入；若日志或终端状态共享索引，
普通自旋锁还会造成“Thread 持锁，被 IRQ 打断，IRQ 等待同一锁”的永久自锁。

\code{IrqSaveSpinLock} 先保存 IF、关闭本地可屏蔽中断，再取得原子锁；释放后
恢复调用前 IF。目标对象注入真实 \code{DisableInterrupts/RestoreInterrupts}，
主机逻辑测试注入无操作函数。这样同一 \code{Terminal} 实现可以在 hosted
测试中运行，却不会把无操作保护带到目标 Kernel。

\begin{contractbox}
锁内只能执行有界内存操作和立即设备写，不能阻塞、sleep 或等待另一个 Thread。
未来异步输出必须把请求移交到锁外队列，而不是扩大当前临界区。
\end{contractbox}

\section{控制终端怎样决定前台进程组}
PID1 的初始身份为：

\[
PID=PGID=SID=1,\qquad session\_leader=true
\]

fork child 继承 SID/PGID，获得独立 PID，并清除 session leader。目标可建立
以自身 PID 为 PGID 的新组，或加入同 SID 已存在的组。跨 session 移动、
移动 session leader 或加入不存在的组都拒绝。

\code{JobControlManager} 与 ProcessTree 使用相同内部 process index，公开
接口仍查找 PID。槽位可以回收，PID 不可倒退复用。完整校验扫描每个 active
记录，证明 PID 唯一、PGID 在同 SID 内有成员、session leader 满足三号相等，
inactive 槽全零且统计数等于实际数。

进程树和 job-control 表分离十分重要：

\begin{itemize}
  \item ProcessTree 的 parent 决定 wait 权限和孤儿重归属；
  \item JobControl 的 PGID 决定终端信号广播；
  \item SID 决定谁有资格修改同一控制终端；
  \item scheduler state 决定谁当前能获得 CPU。
\end{itemize}

把这些字段压进一个枚举会使“停止但仍是同一组成员”“Zombie 仍需被父进程
收集”等合法组合无法表达。

\topic{后台进程为什么不能读取控制终端}
\code{Terminal} 记录 terminal id、controlling SID、session leader PID 和
foreground PGID。只有：

\[
caller.SID=tty.SID\quad\land\quad caller.PGID=tty.foreground.PGID
\]

的进程可以读取。后台组收到固定错误
\code{BACKGROUND\_TERMINAL\_READ=-55}，输入环完全不前进。立即返回比盲目
睡眠更清楚：输入到达不能改变后台身份，只有 Shell 的前台切换才可能满足条件。
完整 Unix 可发送 SIGTTIN；当前先返回专用错误。

输出目前允许后台写。这符合“后台构建仍可输出日志”的基础需求，但可能与 Shell
提示符交错。当前还没有终端显示重绘、SIGTTOU 和 termios \code{TOSTOP}。

\topic{/dev/console 为何进入 VFS}
标准 fd 曾以 ConsoleInput/ConsoleOutput 特殊 kind 直接指向全局对象，用户
不能通过路径观察设备。这里增加最小只读设备后端：

\begin{lstlisting}
VFS
  mount /
  mount /tmp
  mount /dev
    `-- console : CharacterDevice
\end{lstlisting}

VFS 继续负责逐组件路径、mount identity、vnode 引用、stat 与 readdir；
\code{FileDescriptionKind::TerminalDevice} 再把 read/write 转给
\code{Terminal}。字符设备没有普通文件 offset，truncate 返回 Unsupported。
用户只持有 fd，不获得内核对象地址。

当前后端不是完整 devfs：节点在启动期固定，不能动态注册设备，也没有 uid/gid
权限或热插拔。保留 VFS 接口后，后面的 devfs 可以替换后端而不重写 Shell I/O。

\topic{Ctrl-C 和 Ctrl-Z 如何选择目标}
TTY 不保存 PID 列表，只保存 foreground PGID：

\begin{lstlisting}
Ctrl-C -> SendProcessGroupSignal(foreground_pgid, SIGINT)
Ctrl-Z -> SendProcessGroupSignal(foreground_pgid, SIGTSTP)
\end{lstlisting}

SIGTSTP 和 SIGSTOP 的 Default 动作为停止，SIGCONT 的 Default 动作为继续。
STOP 与 Kill 一样不可屏蔽、不可忽略、不可安装 handler；TSTP 可由交互程序
选择 handler。Shell 自己安装 INT/TSTP handler，以免提示符时被默认终止；
fork child 在 exec 前恢复默认动作，所以前台程序仍按普通规则停止或退出。

TTY 不在 IRQ 中遍历全部 Process 并完成上下文切换。IRQ 只产出具名动作，
ProcessRuntime 在既有安全边界调用 SignalManager、ProcessTree 与 scheduler，
把复杂生命周期留在可统一验证的位置。

\topic{Stopped 为什么必须是独立生命周期状态}
停止不是退出，也不是普通 Blocked：

\begin{table}[htbp]
\centering
\small
\begin{osgridtabularx}{\textwidth}{l X X}
状态 & 等待什么 & 资源语义\\ \hline
Blocked & I/O、condition、deadline 等具体条件 & 条件满足后自动 Ready\\ \hline
Stopped & SIGCONT 或作业控制动作 & 全部地址空间、fd 与 CPU 现场保留\\ \hline
Zombie & parent wait 收集 & 执行资源已释放，只保留身份和退出事实\\ \hline
\end{osgridtabularx}
\caption{Blocked、Stopped 与 Zombie 不能合并}
\end{table}

默认停止路径保存 176 字节通用 UserContext 和 512 字节 FXSAVE 现场，把
Process/Thread 标为 Stopped，并从可调度选择中排除。SIGCONT 重新发布原
Thread。它不创建新程序入口，继续的是被停止时的旧指令流。

run queue 可能仍含停止进程以前发布的节点，所以 PopReady 必须跳过它们；
“是否存在可运行 Thread”也必须调用 \code{HasSchedulableReadyThread}，
不能只检查队列非空。否则全是 stopped 节点时 CPU 会在空调度中忙转。

\topic{事件式 wait 如何保留三种事实}
传统 \code{wait} 只在孩子退出时返回，并立即收集 Zombie。Shell 还必须知道
孩子何时停止或继续。公共 56 字节 \code{ProcessWaitEventResult} 保存 PID、
parent PID、event type、termination reason、exit code、exception vector
和 signal number；flags 选择 Exited、Stopped、Continued 与 NoHang。

\begin{pdfdiagram}[node distance=12mm and 18mm]
  \node[os state] (alive1) {Alive};
  \node[os state,right=of alive1] (stopped) {Stopped};
  \node[os state,right=of stopped] (alive2) {Alive};
  \node[os state,right=of alive2] (zombie) {Zombie};
  \draw[os arrow] (alive1) -- node[above,font=\scriptsize]{TSTP/STOP} (stopped);
  \draw[os arrow] (stopped) -- node[above,font=\scriptsize]{CONT} (alive2);
  \draw[os arrow] (alive2) -- node[above,font=\scriptsize]{exit/signal} (zombie);
  \node[os block,below=of stopped] (events) {parent observes\\Stopped, Continued, Exited};
  \draw[os arrow] (stopped) -- (events);
  \draw[os arrow] (alive2) -- (events);
  \draw[os arrow] (zombie) -- (events);
\end{pdfdiagram}
\diagramnote{前两个事件只清 pending；Exited 才回收 Zombie 和槽位。}

孩子可能在父进程重新获得 CPU 前完成 stop、continue 和 exit。ProcessTree
分别保存 stopped/continued pending，exit 不覆盖前两个事实；父进程按顺序
消费。统计同时记录 produced 与 observed，终态必须相等。

\section{Shell 怎样建立和回收前后台作业}
对于 \code{echo data | cat | wc}：

\begin{enumerate}
  \item 在 fork 前取得空闲作业槽并建立两条动态 Pipe；
  \item 第一个 child PID 成为 PGID，后续 child 加入同组；
  \item child 接线 stdin/stdout，关闭多余端点，恢复默认信号，再 exec；
  \item parent 关闭自己的 Pipe 端点并完成每个 child 的 PGID 设置；
  \item 前台命令把 TTY foreground 切到作业 PGID；
  \item Shell 用事件式 wait 汇总全部成员；
  \item 停止或全部退出后，无条件把 TTY foreground 恢复为 Shell PGID。
\end{enumerate}

父子双方调用 setpgid 是经典竞态收缩：child 先运行也能在 exec 前入组，parent
先运行也能在交接终端前建立组。但极短 child 仍可能在 parent 调用前退出。

\topic{控制信号为何必须跨过 exec}
还有一个更短的窗口：parent 已把 TTY 交给新 PGID，而 child 尚未完成 exec。
此时 Ctrl-Z 已经成为 pending 异步事实。exec 可以把旧 Handler 重置为
Default、清除活动 handler frame 和被删除兄弟 Thread 的状态，却不能清掉
Process pending 或 survivor pending；否则替换映像会把用户已经按下的控制键
抹去。QEMU 分步注入把 Ctrl-Z 放进这一窗口，新映像必须在第一次返回用户态前
兑现默认停止。这项测试修正了早期“exec 清空全部 pending”的错误契约。

\topic{为何 Zombie 仍需保留进程组身份}
最初实现在线程退出后立即删除进程级 SignalManager 状态。真实 QEMU 中
\code{unknown} child 以 127 极快退出，parent 随后设置 PGID 得到不存在，
Shell 把它误报为普通执行失败，\code{UNKNOWN\_COMMAND\_REJECTED} 偶发丢失。

添加 sleep 只能改变概率。正确生命周期是：

\begin{lstlisting}
Process exit:
  remove Thread signal state
  keep Process signal/group identity
  mark Zombie

Parent wait Exited:
  reap scheduler slot
  remove Process signal identity
  remove JobControl identity
  collect ProcessTree entry
\end{lstlisting}

Zombie 仍不执行 handler，因为已经没有活动 Thread；保留的少量固定元数据只
作为 PID/PGID 锚点。管线组长即使先退出，后续成员仍能加入相同组。这个修复
把“进程存在多久”从实现细节提升为作业控制契约。

\topic{作业表怎样归约成员状态}
Shell 使用 16 项固定表。每项保存 job id、PGID、最多 16 个 PID、成员状态和
最后 stage 退出码。作业状态按成员归约：

\[
\begin{aligned}
  &\text{全部 Exited} &&\Rightarrow Done\\
  &\text{所有未退出成员均 Stopped} &&\Rightarrow Stopped\\
  &\text{其他} &&\Rightarrow Running
\end{aligned}
\]

\code{command \&} 只允许 \code{\&} 位于整行最后。中间 \code{\&} 返回
BackgroundOperatorNotLast，且在 fork/pipe 前失败。后台作业建立后立即返回
提示符；\code{jobs} 以 NoHang 收割事件；\code{bg} 发送 CONT 但不交出 TTY；
\code{fg} 先交出 TTY，再发送 CONT 并阻塞等待。

\begin{failurebox}
作业表满、pipe/fork 失败、PGID 失败或终端交接失败时，必须关闭已建 fd、终止
已建组、收集可见事件并恢复 Shell 的 foreground PGID。只打印错误而不回滚会
在阶段末留下 Pipe、Zombie 或错误终端所有者。
\end{failurebox}

\topic{Ctrl-Z、fg 与 Ctrl-C 的逐步时序}\par
\begin{lstlisting}
QMP: ctrl-z
  i8042 IRQ1 -> decoder 0x1A
  Terminal -> StopForeground
  SIGTSTP -> foreground PGID
  child Running -> Stopped
  Shell wait observes Stopped
  TTY foreground -> Shell PGID

Shell: fg
  TTY foreground -> job PGID
  SIGCONT -> job PGID
  child Stopped -> Ready
  Shell observes Continued

QMP: ctrl-c
  Terminal -> InterruptForeground
  SIGINT -> job PGID
  child -> Zombie
  Shell observes Exited and reaps
  TTY foreground -> Shell PGID
\end{lstlisting}

QEMU runner 不用固定 sleep 猜测 guest 速度。它为每个输入命令等待新的
\code{COMMAND\_COMPLETE}，控制键前等待新的
\code{FOREGROUND\_JOB\_WAITING}。单步 20 秒无进展立即失败并设置协议完成
事件，finally 终止、wait 并删除 QMP socket；不会再静默等满 120 秒。

\topic{热路径不逐事件打印，只保留事务标记}
扫描码、普通字符和 scheduler 扫描都是热路径。逐事件串口日志会改变时序，
甚至让日志本身成为竞态放大器。低频用户 marker 只描述可作为协议屏障的事务：

\begin{lstlisting}
JOB_CONTROL_READY
FOREGROUND_JOB_WAITING
FOREGROUND_JOB_STOPPED
BACKGROUND_JOB_STARTED
COMMAND_COMPLETE
\end{lstlisting}

全部进程结束后，Kernel 一次输出 TTY、scheduler、signal、ProcessTree 和
JobControl 汇总计数。重要终态包括：

\[
stops=continues,\qquad
stopped\_events=observed\_stopped\_events
\]
\[
continued\_events=observed\_continued\_events,\qquad
output\_queued=output\_written
\]

Stopped、Zombie、active job process、active session/group、editing bytes、
buffered bytes 和 pending output 必须全为零。即使用户 marker 正确，只要
这些计数还有一个非零，整机仍输出 \code{USER\_RESULT\_INVALID}。

\topic{四层测试分别会抓到什么错误}\par
\begin{table}[htbp]
\centering
\small
\begin{osgridtabularx}{\textwidth}{l X}
层次 & 检查内容\\ \hline
单元 & Terminal 行规程/守恒、JobControl 身份、字符 vnode、Ctrl 扫描码\\ \hline
集成 & 同一 child 的 Stopped→Continued→Exited 与 scheduler 状态\\ \hline
随机 & 64 槽、固定种子 100000 步 PGID 迁移，每步完整 Validate\\ \hline
系统 & QMP 真按键、IRQ1、Ring 3 cat、jobs/fg/bg/\&、终态资源守恒\\ \hline
\end{osgridtabularx}
\caption{单元、集成、随机和整机测试覆盖不同故障范围}
\end{table}

functional QEMU 精确观察两次 stop、两次 continue、一次 interrupt，且
Stopped/Continued produced 与 observed 各为两次。后台读取拒绝至少一次。
64 MiB bootstrap 保持兼容，256 MiB 执行完整功能，64 GiB 继续验证同一 ABI
在更大容量配置下运行；QEMU 仍只使用自研 ROM 启动。

\topic{当前终端范围}
当前只有一个控制终端和 canonical 基线，没有 termios/raw mode、多个虚拟
终端或 pty、窗口尺寸、SIGWINCH、SIGHUP、SIGTTIN/SIGTTOU、孤儿进程组规则
与完整 POSIX job spec。Shell 作业表固定 16 项；达到上限或遇到未支持行为时，
系统返回错误。

终端与第 26 章的磁盘请求复用 irq-save 临界区、WaitQueue 单赢家、deadline
和超时处理。两类设备都只在 IRQ 中提交有界结果，Thread 在安全边界恢复；
失败结果保持为失败，不会被唤醒动作改写成成功。
